% Generated draft: requires mathematical and source review.
\documentclass[11pt]{article}
\usepackage[margin=25mm]{geometry}
\usepackage{fontspec}
\setmainfont{Latin Modern Roman}
\usepackage{amsmath,amssymb,mathtools,unicode-math}
\setmathfont{Latin Modern Math}
\usepackage{enumitem,needspace,xurl,hyperref}
\setlength{\parindent}{0pt}
\setlength{\parskip}{5pt}
\setlength{\emergencystretch}{3em}
\newcommand{\N}{\mathbb{N}}
\newcommand{\Z}{\mathbb{Z}}
\newcommand{\Q}{\mathbb{Q}}
\newcommand{\R}{\mathbb{R}}
\newcommand{\C}{\mathbb{C}}
\newcommand{\F}{\mathbb{F}}
\newcommand{\A}{\mathbb{A}}
\newcommand{\PP}{\mathbb{P}}
\newcommand{\E}{\mathbb{E}}
\newcommand{\eps}{\varepsilon}
\newcommand{\dd}{\,\mathrm{d}}
\begin{document}
\section*{UnsolvedMath ID: 30001881}
\textbf{Source problem code:} catalogue entry
\subsection{Definitions and mathematical statement}
Let \(v\) be a positive integer. A Steiner triple system of order \(v\), denoted \(\operatorname{STS}(v)\), is a pair \((V,\mathcal B)\), where \(|V|=v\), \(\mathcal B\subseteq\binom{V}{3}\), and every two-element subset of \(V\) belongs to exactly one member of \(\mathcal B\). Members of \(\mathcal B\) are called blocks. Use the usual admissibility condition \(v\equiv1\) or \(3\pmod 6\).

For an integer \(s\geq4\), a standard precise interpretation of \(s\)-sparsity is
\[
 \forall i\in\Z,\quad 4\leq i\leq s\ \Longrightarrow\ \forall\mathcal F\subseteq\mathcal B,\quad
 |\mathcal F|=i\ \Longrightarrow\ \left|\bigcup_{B\in\mathcal F}B\right|\geq i+3.
\]
Thus no \(i\) distinct blocks are supported on at most \(i+2\) points.

Under this explicitly stated interpretation, determine the set
\[
 \left\{(s,v)\in\Z_{\geq4}\times\Z_{\geq1}:\ v\equiv1,3\pmod6,\ \text{there exists an }s\text{-sparse }\operatorname{STS}(v)\right\}.
\]
The question concerns each admissible order, rather than only existence for sufficiently large orders or for infinite families.

The supplied phrase ``no \(i\) blocks on \(i+2\) points'' does not explicitly distinguish support of size exactly \(i+2\) from support of size at most \(i+2\). The exact-support reading would instead require \(\left|\bigcup_{B\in\mathcal F}B\right|\neq i+2\). A reading in terms of containment in an \((i+2)\)-element subset of \(V\) agrees with the at-most reading when \(i+2\leq v\), but needs separate attention otherwise. The intended convention must therefore be confirmed, especially because no relation between \(s\) and \(v\) is supplied. The source also gives no lower bound for \(s\); the displayed question treats the nonvacuous range \(s\geq4\). For integer \(s<4\), the displayed sparsity condition is vacuous. The positive-order convention includes the degenerate system of order \(1\); the source does not explicitly address its inclusion.
\subsection{Short English statement}
Determine exactly which sparsity levels \(s\) and admissible orders \(v\) permit a Steiner triple system in which every collection of \(i\) distinct blocks, for \(4\leq i\leq s\), uses at least \(i+3\) points. This uses an explicit at-most-support interpretation of the source wording; the intended support convention and small-parameter boundaries require review.
\subsection{Sources}
\begin{itemize}
\item Local UnsolvedMath registry, \texttt{data/unsolvedmath/problems.json}, permanent ID \(30001881\), title ``Existence of Sparse Steiner Triple Systems,'' exact source problem code \texttt{catalogue entry}; supplied complete source record from the Oberwolfach Reports Open Problems collection. No individual proposer is supplied; the proposed year is recorded as 2011.
\item Supplied report citation: \emph{Discrete Geometry (2012)}, Oberwolfach Reports, DOI 10.4171/owr/2011/44, pp.~2459--2548. Original OWR report: \url{https://doi.org/10.4171/owr/2011/44}. The supplied citation labels the report 2012, while the DOI and proposed-year metadata contain 2011; this discrepancy has not been resolved.
\item Supplied literature reference: \emph{On 6-sparse Steiner triple systems}, \url{https://doi.org/10.1016/j.jcta.2006.04.003}. Authors and further bibliographic details are not supplied.
\item Literature assessment dated exactly \texttt{2026-08-21}: the supplied curation labels the problem partially solved and reports existence for several sparsity levels and infinite families, including 6-sparse systems, while describing the general existence question and higher levels as unresolved. This is supplied dated curation, not an independently verified literature assessment; the record explicitly requests independent verification of time-sensitive status and cited claims.
\item The supplied statement-status label is \texttt{corrected\_verified}. Its accompanying verification note states only that TeX notation was enclosed in mathematical delimiters without changing mathematical content; this does not establish mathematical or bibliographic verification of this draft.
\end{itemize}
\end{document}
